% TOP_PROBLEM: {"id": 30006690, "title": "Modularity of Elliptic Curves over Number Fields (End_K(E)=Z)", "category_id": 1, "category": {"id": 1, "name": "number_theory", "display_name": "Number Theory", "description": "Properties of integers, prime numbers, Diophantine equations.", "slug": "number-theory", "order_index": 1, "created_at": "2026-07-31T15:26:25.670Z"}, "status": "open"}
% FIRST_POSTED: 2026-09-27
% !TeX program = lualatex
% Standalone notebook; compile twice: lualatex 30006690.tex
% Original section preserved verbatim except for marked insertion blocks.
\documentclass[11pt,a4paper]{article}
\usepackage[margin=24mm,headheight=14pt]{geometry}
\usepackage{fontspec}
\setmainfont{Latin Modern Roman}
\setsansfont{Latin Modern Sans}
\setmonofont{Latin Modern Mono}
\usepackage{amsmath,amssymb,mathtools}
\usepackage{unicode-math}
\setmathfont{Latin Modern Math}
\usepackage{microtype}
\usepackage{enumitem,xurl,needspace}
\usepackage[dvipsnames]{xcolor}
\usepackage{hyperref}
\hypersetup{unicode=true,colorlinks=true,linkcolor=MidnightBlue,urlcolor=MidnightBlue,
 pdftitle={Research notebook: Problem 30006690},
 pdfauthor={Open Problems Math research notebook},bookmarksnumbered=true}
\usepackage{bookmark,titlesec,fancyhdr}
\titleformat{\section}{\Large\bfseries\raggedright}{\thesection}{0.7em}{}
\pagestyle{fancy}\fancyhf{}
\fancyhead[L]{\small\sffamily Open Problems Math\enspace|\enspace Research notebook}
\fancyhead[R]{\small\sffamily Problem 30006690}
\fancyfoot[L]{\small\sffamily 27 September 2026}
\fancyfoot[R]{\small\sffamily \thepage}
\setlength{\parindent}{0pt}
\setlength{\parskip}{5pt plus 1pt}
\setlength{\emergencystretch}{3em}
\setcounter{secnumdepth}{1}
\setcounter{tocdepth}{1}
\setlist[itemize]{leftmargin=3em,itemsep=5pt,topsep=4pt,parsep=0pt}
\allowdisplaybreaks
\newcommand{\N}{\mathbb{N}}
\newcommand{\Z}{\mathbb{Z}}
\newcommand{\Q}{\mathbb{Q}}
\newcommand{\R}{\mathbb{R}}
\newcommand{\C}{\mathbb{C}}
\newcommand{\F}{\mathbb{F}}
\newcommand{\A}{\mathbb{A}}
\newcommand{\PP}{\mathbb P}
\newcommand{\E}{\mathbb{E}}
\newcommand{\eps}{\varepsilon}
\newcommand{\dd}{\,\mathrm d}
\newcommand{\sref}[2]{\hyperlink{source:#1:#2}{[S#2]}}
\newcommand{\eref}[2]{\hyperlink{extra:#1:#2}{[E#2]}}
\newif\ifshowcatalogue
\showcataloguetrue
\newcommand{\cataloguescope}[1]{\ifshowcatalogue
 \paragraph{Exact scope in the supplied catalogue.}
 \begin{quote}\small #1\end{quote}\fi}
\numberwithin{equation}{section}
\begin{document}
\setcounter{section}{0}
% BEGIN PRESERVED SECTION
% ================================================================
% problemId: 30006690
% Canonical problem ID: 30006690
% exactTarget SHA-256: b671600b2c42b41472f8a1216b710e32313f1dbf5a2a99c7916c13c5acaff8cc
\clearpage
\section*{\texorpdfstring{Modularity of Elliptic Curves over Number Fields (End\_K(E)=Z)}{Modularity of Elliptic Curves over Number Fields (End\_K(E)=Z)}}\label{30006690}
\subsection{Definitions and mathematical statement}
For $E/K$ with $\operatorname{End}_K(E)=\Z$, let $(r_v,N_v)$ be its finite-place Weil--Deligne representations. Use Thorne's geometric-Frobenius convention to form
\[
 \pi(E)=\bigotimes_{v\nmid\infty}'
 \operatorname{rec}_{K_v}^{-1}(r_v\otimes|\cdot|^{1/2},N_v),\qquad
 A=\varinjlim_U H^*(Y_U,\Q).
\]
Here $Y_U$ are the adelic locally symmetric spaces in the source, and the transition maps give the finite-adelic group action. The assertion is the existence of a $\operatorname{GL}_2(\A_K^\infty)$-equivariant injection $\pi(E)\hookrightarrow A\otimes_{\Q}\C$. The condition concerns endomorphisms \emph{over $K$}; replacing it by $\operatorname{End}_{\bar K}(E)=\Z$ would narrow the target.
\par\smallskip\textit{Formulation and concept references:} \sref{30006690}{1}.
\cataloguescope{For every number field K and every elliptic curve E/K with End\_K(E)=Z, prove that pi(E) admits a GL\_2(A\_K\textasciicircum{}infinity)-equivariant injection into A tensor\_Q C, exactly as in Thorne's Conjecture 2.6. Here A is the direct limit over open compact U of H\textasciicircum{}*(Y\_U,Q) for the adelic locally symmetric spaces defined on page 647, and pi(E) is the restricted tensor product at all finite places v of rec\_\{K\_v\}\textasciicircum{}\{-1\}(r\_v tensor |.|\textasciicircum{}\{1/2\},N\_v), with (r\_v,N\_v) attached to E in the source's geometric-Frobenius normalization on pages 646 and 649.}
\subsection{Short English statement}
Does every elliptic curve in the specified number-field class occur in the automorphic cohomology predicted by Thorne's modularity formulation?
% BEGIN RESEARCH 39
\subsection{Research attempt}

\paragraph{Current frontier.}
The incorporated formulation was read in Thorne, pp.~646--650, especially
Definition~2.3 and Conjecture~2.6. It is a statement about an injection of the
entire finite-adelic representation, not just equality of good-prime traces;
Thorne explicitly distinguishes it from his Conjecture~2.5. Geometric
Frobenius and the twist $|\cdot|^{1/2}$ are retained. His Section~4 explains
residual modularity and modularity lifting. \sref{30006690}{1}.
Caraiani--Newton establish modularity over infinitely many imaginary quadratic
fields, including $\Q(\sqrt{-d})$ for $d=1,2,3,5$, under their precise
$X_0(15)$ criterion; their more general CM-field results retain image
hypotheses. This does not quantify over every number field. \eref{30006690}{1}.
Dummigan--Tiwari's August 2026 version proves an implication from arithmetic
hyperbolic uniformization to geometric modularity for totally real fields of
odd degree, not existence of that uniformization for every curve.
\sref{30006690}{4}, \eref{30006690}{2}. These source statements were inspected; their
complete proofs were not independently reverified here.

\paragraph{Attack route.}
Potential automorphy followed by descent would need a descent theorem over the
original field, not an averaging operation on cohomology: changing the base
number field is not the same as passing to a finite covering level $Y_V\to Y_U$.
A second route is to build the exact characteristic-zero eigensystem from
congruences of increasing precision. The attractive feature is that a fixed
level has finite-dimensional cohomology; the difficulty is that neither
bounded level, compatible precision, nor survival beyond torsion comes for
free. I investigate this second route and isolate exactly what compactness
can and cannot supply.

\paragraph{Attempt.}
Let $\mathcal O$ be the integer ring of a finite extension of $\Q_p$, with
uniformizer $\varpi$ and finite residue field. Let $M$ be a finite free
$\mathcal O$-module, let $\mathcal T$ be any family of endomorphisms of $M$,
and prescribe scalars $a_T\in\mathcal O$ for $T\in\mathcal T$. No simultaneous
diagonalization is assumed.

\textbf{Lemma 39.1 (fixed-level, all-precision lifting).}
Suppose that for every $r\ge1$ there is a vector $u_r\in M/\varpi^rM$ whose
reduction modulo $\varpi$ is nonzero and such that
\[
 (T-a_T)u_r=0\quad\text{in }M/\varpi^rM\qquad(T\in\mathcal T).
\]
Then there is $u\in M\setminus\varpi M$ satisfying $(T-a_T)u=0$ for every
$T\in\mathcal T$. In particular the prescribed simultaneous eigenspace in
$M\otimes_{\mathcal O}\operatorname{Frac}(\mathcal O)$ is nonzero.

\emph{Proof.}
At level $r$ take the finite set of all primitive simultaneous solutions.
Reduction maps a solution to a primitive solution at every smaller positive
level. Join each solution to its reduction at the preceding level, adjoining
a root at level zero. The resulting finitely branching tree has vertices at
arbitrarily large depths. To construct an infinite path, choose at each step
a child with descendants at arbitrarily large depths; one exists because
there are only finitely many children. This gives compatible $v_r$ without
assuming that the originally chosen $u_r$ were compatible. Completeness gives
$u\in\varprojlim M/\varpi^rM=M$. For each fixed $T$, $(T-a_T)u$ belongs to
every $\varpi^rM$, whose intersection is zero by freeness. The first reduction
is nonzero, so $u$ is primitive and remains nonzero after inverting $\varpi$.
\hfill$\square$

For a possible application, fix one compact open level $U$ and degree $q$,
and take the torsion-free quotient of an integral cohomology lattice, completed
at $p$, as $M$. At good places set $a_{T_v}=a_v(E)$ in the convention of the
source. The lemma shows that \emph{primitive realizations of this exact system
at every precision, in that one lattice}, would give a characteristic-zero
system. It gives no way to manufacture those realizations from a single
residual modularity statement. Even obtaining the trace eigensystem would
only reach the weaker stage identified by Thorne; all finite local factors,
including their monodromy, still need to be identified before Conjecture~2.6
has been addressed. \sref{30006690}{1}, pp.~648--650.

The torsion issue can be exposed before using any automorphic theorem. For a
bounded complex $C^\bullet$ of finite free $\mathcal O$-modules, the long exact
sequence associated to multiplication by $\varpi^r$ gives
\[
0\longrightarrow H^q(C)/\varpi^rH^q(C)
\longrightarrow H^q(C\otimes\mathcal O/\varpi^r)
\longrightarrow H^{q+1}(C)[\varpi^r]\longrightarrow0.
\]
Here exactness follows by taking the kernel and cokernel of multiplication in
the adjacent terms of that long exact sequence. Therefore residual cohomology
can consist entirely of classes arising from torsion in the next degree.
For the two-term complex $\mathcal O\xrightarrow{\varpi}\mathcal O$ in degrees
$q,q+1$, the rationalized complex is acyclic whereas its reduction modulo
$\varpi$ has one-dimensional cohomology in both degrees. A nonzero mod-$p$
class alone thus does not establish a nonzero rational class.

There is also a prescribed-eigenvalue obstruction, even with no torsion. Take
$M=\Z_p$, $T=0$, and prescribe $a_T=p^s$ for a positive integer $s$. A primitive
solution exists modulo $p^r$ precisely for $r\le s$, and there is no nonzero
solution over $\Q_p$. Indeed $p^su\equiv0\pmod{p^r}$ for a unit $u$ holds
exactly when $r\le s$. Thus an arbitrarily impressive \emph{finite} precision
computation cannot replace the all-precision quantifier. The accompanying
script checks this equivalence for small primes and precisions by exact
integer arithmetic, as a test of the example rather than of modularity.

\paragraph{Stress test.}
The proof uses finite residue fields, fixed finite rank, fixed level, and
primitivity in a torsion-free lattice. Letting $U=U_r$ vary destroys the single
compact tree. Permitting $u_r$ to become more divisible allows all normalized
mass to vanish. Working in torsion cohomology can yield an inverse-limit vector
that dies in characteristic zero. Each would invalidate a proposed application.
Furthermore, potential automorphy over an extension does not identify a
$\operatorname{GL}_2(\A_K^\infty)$-subrepresentation over $K$ by elementary
Galois averaging. A full proof must preserve the ramified Weil--Deligne
operators $N_v$, not merely Frobenius semisimplifications at good places.
The hypothesis $\operatorname{End}_K(E)=\Z$ is unchanged: curves acquiring
complex multiplication only after extending $K$ have not been removed.

One possible next lemma would assert, for every curve in the entry, a single
level and an integral local-type condition admitting primitive eigensystems
modulo every $p^r$, with characteristic-zero local-global compatibility at all
finite places. That is a substantial automorphy assertion, not a consequence
of Lemma~39.1. It is intentionally left unproved rather than hidden in a claim
that compatible lifts are automatic.

\Needspace{8\baselineskip}
\paragraph{Outcome.}
\textbf{Outcome: Exploratory approach with unresolved gap.}

A fixed-level lifting lemma and two explicit obstructions to naive residual
lifting have been proved. Historical novelty of this elementary lemma is not
claimed. The attempt has not constructed the required primitive systems, nor
identified all finite local factors; consequently it has not supplied the
requested injection. The useful reduction is that compactness itself is not
the obstacle once those exact hypotheses are met: the obstacle is producing
them while preventing torsion, level growth, and loss of local information.

% END RESEARCH 39

\subsection{Sources}
\begin{itemize}\raggedright
\item[\textnormal{[S1]}]\hypertarget{source:30006690:1}{} J. A. Thorne, Elliptic curves and modularity, 8ECM proceedings, EMS Press, 2023, DOI 10.4171/8ECM/12.
\url{https://ems.press/content/book-chapter-files/28276}.
{\small\textit{Catalogue formulation source.}}
\item[\textnormal{[S2]}]\hypertarget{source:30006690:2}{} Elliptic curves and modularity.
\url{https://ems.press/content/book-chapter-files/28276?nt=1}.
{\small\textit{Catalogue research/status reference; not a proof endorsement.}}
\item[\textnormal{[S3]}]\hypertarget{source:30006690:3}{} Modularity of elliptic curves over cyclotomic Z\_p-extensions of real quadratic fields.
\url{https://arxiv.org/html/2206.12860v1}.
{\small\textit{Catalogue research/status reference; not a proof endorsement.}}
\item[\textnormal{[S4]}]\hypertarget{source:30006690:4}{} Modular elliptic curves and hyperbolic uniformization.
\url{https://arxiv.org/abs/2607.03830}.
{\small\textit{Catalogue research/status reference; not a proof endorsement.}}
% BEGIN ADDED REFERENCES 39
\item[\textnormal{[E1]}]\hypertarget{extra:30006690:1}{} Ana Caraiani and James Newton,
\emph{On the modularity of elliptic curves over imaginary quadratic fields},
arXiv:2301.10509, abstract and main theorem statements.
\url{https://arxiv.org/abs/2301.10509}.
{\small\textit{Consulted 27 September 2026; the field and residual-image restrictions are retained.}}
\item[\textnormal{[E2]}]\hypertarget{extra:30006690:2}{} Neil Dummigan and Devendra Tiwari,
\emph{Modular elliptic curves and hyperbolic uniformization},
arXiv:2607.03830v2, 27 August 2026, abstract and stated implication.
\url{https://arxiv.org/abs/2607.03830v2}.
{\small\textit{Consulted 27 September 2026; conditional uniformization implication, not universal modularity.}}

% END ADDED REFERENCES 39
\end{itemize}

% END PRESERVED SECTION
\end{document}
